\section{Punkt odniesienia: MFW i modelowanie BDP}

Raport Międzynarodowego Funduszu Walutowego \textit{Fiscal Monitor: Tackling Inequality} analizuje BDP jako instrument redystrybucyjny, którego skutki należy oceniać przez koszty fiskalne, efektywność, wpływ na podaż pracy oraz zmianę nierówności \parencite{imf2017fiscal}. Szczególnie istotny jest Aneks 1.3, wykorzystujący dynamiczny stochastyczny model równowagi ogólnej zaadaptowany z pracy Lizarazo, Peralta-Alvy i Puya \parencite{lizarazo2017}.

W tej strukturze gospodarka jest podzielona na trzy sektory: sektor dóbr podlegających wymianie międzynarodowej, usługi wymagające niskich kwalifikacji oraz usługi wymagające wysokich kwalifikacji. Gospodarstwa domowe różnią się poziomem wykształcenia i są obarczone indywidualnymi szokami produktywności. Rynki pracy są segmentowane, a krajowe rynki kredytowe są niekompletne. Model porównuje wariant odniesienia z wariantami reform podatkowo-transferowych, w tym BDP finansowanym cięciami wydatków, wzrostem VAT lub wzrostem progresywnego PIT.

Wniosek MFW jest wewnętrznie spójny: BDP jest progresywny, ale przy tej samej skali kosztu fiskalnego może być zdominowany dobrobytowo przez instrumenty celowane, takie jak rozszerzenie Earned Income Tax Credit. W takim ujęciu pytanie brzmi: jaki instrument redystrybucyjny maksymalizuje dobrobyt społeczny przy zadanym koszcie i założonych mechanizmach finansowania?

Polska w raporcie MFW pojawia się w innej formule. Aneks 1.6 nie jest pełnym DSGE, lecz analizą częściowej równowagi statycznej (partial static equilibrium) opartą na danych Luxembourg Income Study dla Polski z 2013 r. Koszt brutto BDP skalibrowanego do 25\% mediany dochodu rynkowego po podatkach bezpośrednich na osobę wynosi 4,9\% PKB, redukcja współczynnika Giniego wynosi 0,04, redukcja stopy ubóstwa 6,9 punktu procentowego, a kwota BDP to 2 111 zł na osobę rocznie. Nominalnie odpowiada to około 176 zł na osobę miesięcznie, ale jest to wartość osadzona w danych z 2013 r. Po przybliżonej indeksacji rocznymi wskaźnikami CPI GUS za lata 2014--2025 daje to około 273 zł na osobę miesięcznie w cenach średniorocznych 2025 \parencite{gusCpiAnnual}. Właśnie dlatego nie należy porównywać nominalnych 176 zł bezpośrednio z poziomami siatki scenariuszy \amor{}, które są miesięcznym transferem na agenta reprezentującego gospodarstwo domowe w modelu skalibrowanym do Polski na 2026-04-30.

Most między MFW i \amor{} nie jest więc mapowaniem jeden do jednego. MFW Polska 2013 i \amor{} Polska 2026 odpowiadają na różne pytania, mają różne lata kalibracji i używają różnych jednostek transferu. Oryginalna kotwica fiskalna MFW, czyli 4,9\% PKB w roku kalibracji 2013, jest wartością historyczną dla tamtego ćwiczenia. Po indeksacji CPI kwota transferu wynosi około 273 zł na osobę miesięcznie; przy nominalnym PKB Polski 2025 z Article IV MFW daje to około 3,2\% PKB \parencite{imf2026polandArticleIV}. Tabela~\ref{tab:mfw-scale-bridge} pokazuje także pomocnicze odniesienie tej samej kwoty do wspólnego mianownika modelowego, czyli PKB wariantu bez BDP po 60 miesiącach w \amor{}; daje to około 2,32\% PKB. To nie jest scenariusz MFW ani wynik wariantu \amor{}, lecz przeliczenie pomostowe. Dlatego w dalszej interpretacji nie traktuję żadnego poziomu siatki jako dokładnej repliki MFW. BDP 500~PLN jest najbliższą kotwicą w ujęciu na osobę i kotwicą w porównywalnym mianowniku modelowym, a BDP 1000~PLN pozostaje blisko oryginalnej, historycznej kotwicy 4,9\% PKB z aneksu MFW.

Dla poziomów siatki \amor{} kolumna „Ekwiwalent na osobę” w tabeli~\ref{tab:mfw-scale-bridge} nie jest parametrem wejściowym modelu. Jest liczona jako dodatkowy miesięczny strumień transferów społecznych względem wariantu bez BDP, podzielony przez populację referencyjną 38 mln osób. Dlatego poziom BDP 500~PLN na agenta reprezentującego gospodarstwo domowe odpowiada około 283 PLN/os./mies., a nie 500 PLN/os./mies.

\begin{table}[htbp]
\centering
\scriptsize
\caption{Skala polskiego wariantu MFW względem siatki eksperymentu \amor{}}
\label{tab:mfw-scale-bridge}
\begin{tabularx}{\textwidth}{@{}p{0.20\textwidth}rrrX@{}}
\toprule
Pozycja & Ekwiwalent na osobę & Koszt wg źródła & Koszt porównawczy & Interpretacja \\
 & PLN/os./mies. & \% PKB źródła & \% PKB bez BDP, m60 & \\
\midrule
MFW, Polska, ceny 2013 & 176 & 4,9 & -- & kwota MFW: 2 111 PLN/os./rok; mianownik PKB 2013 \\
MFW, Polska, ceny 2025 & 273 & 3,2 & 2,32 & kwota MFW po CPI; kolumna 4 to tylko wspólny mianownik modelu \\
BDP 500~PLN & 283 & -- & 2,40 & poziom siatki najbliższy MFW po CPI \\
BDP 1000~PLN & 565 & -- & 4,80 & blisko historycznej kotwicy 4,9\% PKB MFW \\
BDP 2000~PLN & 1 130 & -- & 9,59 & lokalny szczyt adopcji AI/hybrid w wariancie centralnym \\
BDP 3000~PLN & 1 700 & -- & 14,42 & wysoki impuls fiskalno-finansowy \\
\bottomrule
\end{tabularx}
\end{table}


Tabela~\ref{tab:mfw-scale-bridge} rozdziela dwa mianowniki: koszt w źródle oraz koszt przeliczony na wspólny mianownik modelowego PKB wariantu bez BDP po 60 miesiącach. Po takim rozdzieleniu widać, że polski wariant ilustracyjny MFW po indeksacji CPI znajduje się blisko najniższego dodatniego poziomu siatki \amor{} w ujęciu na osobę i w modelowym mianowniku PKB. Jednocześnie oryginalny koszt brutto MFW, liczony jako udział w PKB roku kalibracji 2013, pozostaje blisko scenariusza BDP 1000~PLN. Oznacza to, że most MFW--\amor{} ma charakter zakresu, a nie jednego punktu. Poziomy BDP od 1000 PLN do 3000 PLN testują coraz silniejsze impulsy fiskalno-dochodowe. Ich rolą nie jest odtworzenie polskiego ćwiczenia statycznego MFW, lecz przetestowanie progu, przy którym pozytywny impuls popytowy zaczyna ustępować interakcji presji płacowo-rezerwowej i ograniczeń fiskalno-finansowych.

\begin{table}[H]
\centering
\small
\caption{MFW i \amor{} jako komplementarne narzędzia analityczne}
\label{tab:mfw-vs-amor-fati}
\begin{tabularx}{\textwidth}{@{}L{0.20\textwidth}L{0.13\textwidth}>{\raggedright\arraybackslash}X@{}}
\toprule
Wymiar & Narzędzie & Ujęcie \\
\midrule
Pytanie badawcze & MFW & Jaki instrument redystrybucyjny poprawia równość i dobrobyt przy zadanym koszcie fiskalnym? \\
 & \amor{} & Jak BDP przechodzi przez bilanse, koszt kapitału, kurs, banki, inwestycje i decyzje firm o automatyzacji? \\
\addlinespace
Perspektywa czasu & MFW & Porównanie wariantów polityki w równowadze lub w analizie statycznej. \\
 & \amor{} & Ścieżka przejścia w horyzoncie 60 miesięcy. \\
\addlinespace
Jednostka BDP & MFW & Transfer na osobę, w polskim aneksie raportowany rocznie. \\
 & \amor{} & Miesięczny transfer na agenta reprezentującego gospodarstwo domowe, przeliczany następnie na ekwiwalent na osobę. \\
\addlinespace
Finansowanie & MFW & Warianty finansowania przez cięcia wydatków, VAT lub PIT; dla Polski koszt brutto w analizie statycznej. \\
 & \amor{} & Wariant deficytowy względem ścieżki bez BDP w \amor{}, z jawnym przejściem przez wydatki, deficyt i dług. \\
\addlinespace
Wyniki & MFW & Redystrybucja, ubóstwo, Gini, dobrobyt, koszt fiskalny. \\
 & \amor{} & Adopcja AI/hybrid, inwestycje, inflacja, stopa NBP, dług, deficyt, kurs, rachunek bieżący i bilanse banków. \\
\bottomrule
\end{tabularx}
\end{table}


Z perspektywy niniejszej pracy MFW nie jest przeciwnikiem, lecz punktem odniesienia. Raport odpowiada precyzyjnie na pytanie o redystrybucję i dobrobyt społeczny. \amor{} odpowiada na inne pytanie: czy firmy w trakcie przejścia mają finansową i organizacyjną zdolność do automatyzacji. To pytanie wymaga modelu, w którym ścieżka przejścia jest częścią wyniku.
